%=============================================================================
% MÓDULO 03: INTRODUÇÃO (Expansão Adicional)
%=============================================================================
% Frontiers in Public Health — Estrutura de funil com contexto expandido
%   contexto amplo → problema → lacuna → objetivos → contribuição
%=============================================================================

\section{Introdução}

\subsection{Contexto}

A inteligência artificial (IA) e o aprendizado de máquina (AM) estão 
fundamentalmente transformando a prestação de serviços de saúde em múltiplos 
domínios, desde imagem diagnóstica até sistemas de apoio à decisão clínica 
\cite{topol2019}. O rápido avanço das capacidades computacionais, acoplado 
ao crescimento exponencial de prontuários eletrônicos e dados biomédicos, 
criou oportunidades sem precedentes para desenvolver modelos preditivos que 
possam auxiliar os clínicos a fazer diagnósticos mais precisos e oportunos. 
Somente nos Estados Unidos, dispositivos médicos habilitados por IA receberam 
mais de 500 aprovações da Food and Drug Administration, com aplicações 
abrangendo radiologia, cardiologia, patologia e atenção primária.

No cuidado ao diabetes especificamente, os modelos de aprendizado de máquina 
demonstram promessa substancial para detecção precoce, estratificação de 
risco e recomendações de tratamento personalizadas \cite{sd2019, esteva2017}. 
A epidemia global de diabetes, afetando mais de 537 milhões de adultos em 
todo o mundo com projeções indicando aumento para 783 milhões até 2045 
\cite{idf2021}, demanda ferramentas de triagem escaláveis e precisas que 
possam operar efetivamente em populações diversas. Estimativas atuais 
sugerem que aproximadamente um em cada dois adultos com diabetes permanece 
não diagnosticado, representando um desafio significativo de saúde pública 
que ferramentas de triagem baseadas em IA poderiam potencialmente resolver.

As abordagens de aprendizado de máquina para predição de diabetes evoluíram 
significativamente nas últimas duas décadas. Estudos iniciais dependiam de 
métodos estatísticos tradicionais como regressão logística e árvores de 
decisão \cite{smith1988}, enquanto trabalhos mais recentes exploraram 
métodos ensemble \cite{breiman2001, friedman2001}, máquinas de vetores 
de suporte \cite{vapnik1995} e redes neurais profundas. Esses modelos 
mostraram resultados promissores, com precisões relatadas variando de 75\% 
a 95\% dependendo do conjunto de dados, características e algoritmo 
utilizado. No entanto, o foco em métricas de desempenho geral ofuscou 
em grande medida questões críticas sobre equidade e justiça em predições 
algorítmicas.

No entanto, evidências crescentes sugerem que essas tecnologias podem 
perpetuar ou amplificar disparidades existentes de saúde através de viés 
algorítmico \cite{mehrabi2021}. O viés algorítmico em saúde refere-se à 
discriminação sistemática e injusta contra certos grupos, frequentemente 
enraizada em dados históricos que refletem desigualdades sociais 
\cite{barocas2019}. Este fenômeno é particularmente preocupante em contextos 
clínicos, onde predições enviesadas podem impactar diretamente os desfechos 
dos pacientes através de diagnósticos atrasados, tratamentos inadequados 
e alocação desigual de recursos.

As consequências do viés algorítmico em saúde são abrangentes e 
potencialmente devastadoras. Quando um modelo de diagnóstico apresenta 
desempenho sistematicamente inferior para certos grupos demográficos, ele 
cria um sistema de saúde em dois níveis onde as predições algorítmicas são 
confiáveis para populações privilegiadas mas não confiáveis para comunidades 
marginalizadas. Isso pode exacerbar disparidades existentes de saúde, 
erosionar a confiança nos sistemas de saúde e, em última análise, minar a 
promessa de IA para melhorar a saúde para todos.

No diagnóstico de diabetes, modelos enviesados podem produzir predições 
menos precisas para populações sub-representadas, levando a diagnósticos 
atrasados, tratamentos inadequados e, em última análise, piores desfechos 
de saúde \cite{obermeyer2019}. Os mecanismos através dos quais o viés se 
manifesta em modelos de predição de diabetes são multifacetados e incluem 
viés de representação de dados, viés de medição e viés algorítmico. 
Compreender esses mecanismos é essencial para desenvolver estratégias 
eficazes de mitigação.

A epidemia de diabetes afeta desproporcionalmente comunidades 
marginalizadas, criando um contexto particularmente urgente para abordar 
o viés algorítmico. Nos Estados Unidos, a prevalência de diabetes é 60\% 
maior entre adultos negros e 77\% maior entre adultos hispânicos em 
comparação com adultos brancos \cite{cdc2020}. Essas disparidades são 
impulsionadas por determinantes sociais da saúde, incluindo acesso à saúde, 
nutrição e fatores socioeconômicos \cite{hill2022}. Quando os modelos de 
AM são treinados com dados que sub-representam essas populações, eles 
podem falhar em generalizar efetivamente, exacerbando desigualdades 
existentes e criando ciclos de retroalimentação que perpetuam disparidades 
de saúde.

\subsection{Declaração do Problema}

Apesar da crescente conscientização sobre justiça algorítmica, poucos 
estudos avaliaram sistematicamente o viés em modelos de predição de diabetes 
em múltiplas dimensões demográficas simultaneamente. A maioria das pesquisas 
existentes concentra-se em atributos individuais (raça \textit{ou} gênero) 
em vez de examinar efeitos interseccionais onde múltiplas formas de 
desvantagem convergem \cite{crenshaw1989, buolamwini2018}. Essa abordagem 
de atributo único falha em capturar a natureza complexa e acumulativa da 
discriminação experimentada por indivíduos que pertencem a múltiplos grupos 
marginalizados.

Além disso, embora métricas técnicas de justiça tenham sido propostas 
\cite{corbett2018}, suas implicações práticas para a decisão clínica 
permanecem insuficientemente compreendidas. A lacuna entre as definições 
teóricas de justiça e suas aplicações clínicas no mundo real representa 
uma barreira crítica para implementar sistemas de IA equitativos em 
ambientes de saúde.

O arcabouço da interseccionalidade, originalmente articulado por Crenshaw 
\cite{crenshaw1989}, destaca como identidades sobrepostas criam modos 
únicos de discriminação que não podem ser capturados analisando cada 
atributo isoladamente. No contexto da justiça algorítmica, isso significa 
que o viés experimentado por indivíduos na intersecção de múltiplos grupos 
marginalizados (por exemplo, mulheres minoritárias) pode exceder a soma 
dos vieses atribuídos a cada grupo individual. Apesar de sua importância 
na compreensão da discriminação, a análise interseccional continua 
sub-representada na pesquisa sobre justiça de IA em saúde.

\subsection{Objetivos}

Este estudo tem como objetivo:

\begin{enumerate}[leftmargin=*]
    \item Avaliar o desempenho discriminativo de quatro modelos de AM para 
          diagnóstico de diabetes em subgrupos demográficos definidos por 
          gênero e etnia, fornecendo uma avaliação abrangente de como 
          diferentes abordagens algorítmicas se comportam em populações 
          diversas.
    \item Quantificar o viés algorítmico utilizando métricas estabelecidas 
          de justiça, incluindo Odds Equalizados e Paridade Demográfica, 
          para determinar se existem disparidades sistemáticas nas predições 
          dos modelos.
    \item Realizar análise interseccional para identificar disparidades 
          amplificadas na intersecção de raça e gênero, examinando se 
          indivíduos na convergência de múltiplas identidades marginalizadas 
          experimentam desvantagem acumulada.
    \item Fornecer recomendações baseadas em evidências para mitigação de 
          viés em sistemas clínicos de AM implantados no cuidado ao diabetes, 
          oferecendo orientação prática para desenvolvedores e profissionais 
          de saúde.
\end{enumerate}

\subsection{Significado e Contribuição}

Esta pesquisa contribui para a literatura crescente sobre equidade de IA 
em saúde através de:

\begin{itemize}[leftmargin=*]
    \item Um arcabouço de avaliação abrangente para avaliar viés algorítmico 
          que integra desempenho, justiça e análise interseccional, 
          fornecendo um modelo metodológico para pesquisas futuras.
    \item Evidências empíricas de disparidades interseccionais em predição 
          de diabetes, indo além da avaliação de viés de atributo único para 
          capturar a natureza complexa da discriminação.
    \item Orientação prática para desenvolver sistemas de IA clínica mais 
          justos, incluindo recomendações aplicáveis para coleta de dados, 
          desenvolvimento de modelos e monitoramento de implantação que 
          podem ser imediatamente aplicadas em contextos clínicos.
    \item Uma metodologia totalmente reprodutível com código de fonte 
          aberta, facilitando a replicação e extensão pela comunidade de 
          pesquisa e permitindo a validação de descobertas em diferentes 
          contextos.
\end{itemize}

\subsection{Trabalhos Relacionados}

O estudo da justiça algorítmica em saúde ganhou atenção significativa nos 
últimos anos, com vários estudos emblemáticos destacando a prevalência e 
consequências do viés em sistemas de IA clínica. Obermeyer et al. 
\cite{obermeyer2019} demonstraram que um algoritmo de saúde amplamente 
utilizado subestimou sistematicamente as necessidades de saúde de pacientes 
negros, afetando milhões de indivíduos anualmente. Essa descoberta gerou 
debate generalizado sobre as implicações éticas da decisão algorítmica 
em saúde e a necessidade de abordagens conscientes de justiça.

No domínio da predição de diabetes especificamente, vários estudos avaliaram 
o viés entre grupos demográficos. Chen et al. \cite{chen2019} descobriram 
que modelos de predição treinados com prontuários eletrônicos mostraram 
desempenho diferenciado entre grupos raciais, com taxas de falso negativo 
maiores para populações minoritárias. Da mesma forma, Rajkomar et al. 
\cite{rajkomar2018} identificaram disparidades de desempenho em modelos 
clínicos de AM entre subgrupos demográficos em sua revisão sistemática de 
aplicações de IA em saúde.

O arcabouço da interseccionalidade foi aplicado à justiça algorítmica em 
outros domínios, mais notavelmente por Buolamwini e Gebru \cite{buolamwini2018}, 
que demonstraram disparidades significativas de desempenho para mulheres de 
pele escura em sistemas de reconhecimento facial. No entanto, a análise 
interseccional continua rara na pesquisa de IA em saúde, representando uma 
lacuna crítica que este estudo aborda.

Nosso trabalho baseia-se nessas fundações fornecendo uma avaliação 
abrangente de viés algorítmico em predição de diabetes que integra desempenho, 
justiça e análise interseccional. Diferente de estudos anteriores que 
concentraram-se em atributos individuais ou métricas específicas de justiça, 
examinamos múltiplas dimensões simultaneamente para capturar a natureza 
complexa do viés em sistemas de IA clínica.
